% Side-Path Parametric Memory for LLM Agents
% Compile:  pdflatex main && pdflatex main   (from this directory)
% Numbers are generated into generated_tables.tex by paper/build_tables.py;
% this file only ever refers to the emitted macros, never to a hand-copied value.
\documentclass[11pt]{article}

\usepackage[margin=1in]{geometry}
\usepackage{booktabs}
\usepackage{amsmath}
\usepackage{amssymb}
\usepackage[hidelinks]{hyperref}
\usepackage{microtype}
\usepackage{xcolor}

% NOTE: generated_tables.tex is \input *inside* the document body: it contains
% table floats and \newcommand macros, and floats are illegal in the preamble.
% Both work in the body, as long as nothing refers to a macro before this point.

\title{\bf Side-Path Parametric Memory for LLM Agents:\\
Exactly Erasable Slots, Read Without Context}
% Author line: real name + affiliation (user decision, 2026-09-29). Deliberately
% no email: the user's own privacy rule keeps contact details out of public material.
\author{Haotian Cheng \\ \small Nanjing University of Aeronautics and Astronautics}
\date{2026-09-29}

\begin{document}
\maketitle
% Generated by parammem.report.build_results_latex -- do not edit by hand.
% Rebuild with: python paper/build_tables.py

% ---- headline numbers (generated) ----
\newcommand{\tTargetCE}{0.0000}
\newcommand{\tEraseIdentical}{10/10}
\newcommand{\tIsolation}{ok}
\newcommand{\tOracle}{15/16}
\newcommand{\tSum}{0/16}
\newcommand{\tTopOne}{14/16}
\newcommand{\tTopTwo}{4/16}
\newcommand{\tRouterAcc}{0.938}
\newcommand{\tAlwaysWrites}{18}
\newcommand{\tAlwaysWasted}{13}
\newcommand{\tSelfWrites}{5}
\newcommand{\tSelfWasted}{0}
\newcommand{\tSelfMissed}{0}
\newcommand{\tAlwaysKL}{2.59}
\newcommand{\tSelfKL}{0.54}
\newcommand{\tErased}{6/6}
\newcommand{\tErasedTotal}{30/30}
\newcommand{\tPolicies}{5}
\newcommand{\tFifoHot}{1/1}
\newcommand{\tSmartHot}{3/3}
\newcommand{\tFloor}{0.00}
\newcommand{\tContext}{1.00}
\newcommand{\tParam}{1.00}
\newcommand{\tWithholdFloor}{0.17}
\newcommand{\tWithholdMemory}{0.00}
\newcommand{\tScreenKept}{25/30}
\newcommand{\tVirgin}{True}
\newcommand{\tSixOracle}{6/6}
\newcommand{\tRouted}{5/6}

\begin{table}[t]
\centering
\caption{Stage status of this results bundle.}
\label{tab:status}
\begin{tabular}{ll}
\toprule
stage & status \\
\midrule
t1 & ok \\
t2 & ok \\
t3 & ok \\
t4 & ok \\
t5 & ok \\
t6 & ok \\
\bottomrule
\end{tabular}
\end{table}

\begin{table}[t]
\centering
\caption{T1: the sum-based read (all slots active) versus the two controls. The low \texttt{mem\_on} here is the superposition failure that T2 fixes; it is not evidence about whether writing works.}
\label{tab:t1}
\begin{tabular}{lrr}
\toprule
arm & EM & hits/total \\
\midrule
\texttt{prompt\_only} & 0.000 & 0/10 \\
\texttt{mem\_on} & 0.100 & 1/10 \\
\texttt{mem\_shuffle} & 0.000 & 0/10 \\
\texttt{mem\_off} & 0.000 & 0/10 \\
\bottomrule
\end{tabular}
\end{table}

\begin{table}[t]
\centering
\caption{T2: composition rules on the same writes, read with a paraphrased query. Selecting exactly one slot matches the oracle; summing destroys recall; two slots already lose most of it.}
\label{tab:t2}
\begin{tabular}{lrr}
\toprule
arm & EM & hits/total \\
\midrule
\texttt{oracle} & 0.938 & 15/16 \\
\texttt{all} & 0.000 & 0/16 \\
\texttt{top1} & 0.875 & 14/16 \\
\texttt{top2} & 0.250 & 4/16 \\
\texttt{other} & 0.000 & 0/16 \\
\bottomrule
\end{tabular}
\end{table}

\begin{table}[t]
\centering
\caption{T3: write criteria. ``wasted on known'' counts gradient writes spent on facts the model already answers; ``anchor KL'' is the price paid by the frozen backbone.}
\label{tab:t3}
\begin{tabular}{lrrrrrr}
\toprule
criterion & writes & target & wasted & missed & recall & anchor KL \\
\midrule
\texttt{always} & 18 & 5/5 & 13/13 & 0 & 5/5 & 2.59 \\
\texttt{surprise} & 9 & 3/5 & 6/13 & 2 & 3/3 & 1.20 \\
\texttt{selfcheck} & 5 & 5/5 & 0/13 & 0 & 5/5 & 0.54 \\
\texttt{explicit} & 1 & 1/5 & 0/13 & 4 & 1/1 & 0.10 \\
\bottomrule
\end{tabular}
\end{table}

\begin{table}[t]
\centering
\caption{T3: pre-write loss by class. The ranges overlap almost completely, which is why the likelihood-based criterion cannot separate them.}
\label{tab:t3loss}
\begin{tabular}{lrrrr}
\toprule
class & n & min & median & max \\
\midrule
unknown & 5 & 2.64 & 7.72 & 13.92 \\
known & 13 & 3.65 & 6.16 & 10.60 \\
\bottomrule
\end{tabular}
\end{table}

\begin{table}[t]
\centering
\caption{T4: capacity and forgetting. 12 memories stream into 6 slots; the first 3 are accessed repeatedly. Every eviction returns its slot to the virgin state bit-for-bit.}
\label{tab:t4}
\begin{tabular}{lrrrrrrcr}
\toprule
policy & evictions & erased & accessed & untouched & routed & stamp & anchor KL \\
\midrule
\texttt{fifo} & 6 & 6/6 & 1/5 & 1/5 & 5/6 & yes & 0.97 \\
\texttt{lru} & 6 & 6/6 & 3/3 & 3/3 & 5/6 & yes & 1.02 \\
\texttt{lfu} & 6 & 6/6 & 3/3 & 3/3 & 5/6 & yes & 1.30 \\
\texttt{utility} & 6 & 6/6 & 3/3 & 3/3 & 5/6 & yes & 1.02 \\
\texttt{utility\_time} & 6 & 6/6 & 3/3 & 3/3 & 5/6 & yes & 1.02 \\
\bottomrule
\end{tabular}
\end{table}

\begin{table}[t]
\centering
\caption{T5: context inertia, on scenarios whose floor (no premise anywhere) is zero. Storing the premise in the weights drags it into the next topic exactly as much as leaving it in the context, so the hypothesis is refuted.}
\label{tab:t5}
\begin{tabular}{lrrr}
\toprule
arm & inheritance & residue & withholding \\
\midrule
\texttt{no\_memory} & 0.00 & 0.00 & 0.17 \\
\texttt{context\_memory} & 1.00 & 0.42 & 0.00 \\
\texttt{param\_memory} & 1.00 & 0.42 & 0.00 \\
\texttt{param\_off} & 0.00 & 0.00 & 0.17 \\
\bottomrule
\end{tabular}
\end{table}

\begin{table}[t]
\centering
\caption{T6: persistence across two separate processes. The read session's side path is virgin before loading, so the recall it achieves can only come from disk.}
\label{tab:t6}
\begin{tabular}{ll}
\toprule
measurement & value \\
\midrule
oracle recall after restart & 6/6 \\
routed recall (k=1) & 5/6 \\
routing correct & 5/6 \\
side path virgin before load & True \\
erase after load is virgin & True \\
erased memory still recalled & False \\
snapshot size & 18.4 MB (8 slots) \\
\bottomrule
\end{tabular}
\end{table}


\begin{abstract}
Long-term memory in LLM agents lives in the context window, where it competes with the
tokens needed to reason and is indistinguishable from the current working set. We study
an alternative: a capacity-bounded side path of low-rank slots on a frozen backbone,
written by gradient steps, read through the forward pass, and forgotten by erasing a rank
block. On a 0.6B model, writes touch only their own slot; erasure is bit-identical to a
never-written slot across \tErasedTotal{} evictions and a process restart; and the binding
constraint turns out to be read-time composition, since summing the slots destroys recall
($\tSum{}$) while selecting one by query-key similarity recovers it ($\tTopOne{}$ against
an oracle of $\tOracle{}$), with the residual gap equal to the routing error. A label-free
self-check reaches full-write recall using $\tSelfWrites$ writes instead of
$\tAlwaysWrites$, and $\tSelfKL$ nats of backbone drift instead of $\tAlwaysKL$. Two
hypotheses fail, and we report them: a likelihood-threshold surprise criterion measures
wording rather than knowledge, and parametric memory does not reduce context inertia. On
\tLongMemSubset{} LongMemEval single-session questions with the history absent, the
written arm reproduces $\tLongMemOn{}$ of the reference answers at no extra prompt-token
cost, while context delivery answers $\tJudgeContextTarget{}$ to $\tJudgeContextAll{}$ and
pays $\tTokContextTarget{}$ to $\tTokContextAll{}$ tokens. Two checks round it out: at
equal item-gradient-steps the bank recalls $\tBaseSlots{}$ of those thirty memories where
full fine-tuning reaches $\tBaseFull{}$ and one equal-rank adapter $\tBaseAdapter{}$, and
a model-free change of retrieval key lifts 1.7B routing from $\tKeyLastTopLarge{}$ to
$\tKeyBestTopLarge{}$.
\end{abstract}

\section{Introduction}

Two problems motivate this work, and we separate them because our results treat
them very differently.

\paragraph{The context channel is zero-sum.} Without changing weights, the only
place new information can go is the context window, and that channel has three
structural defects. It is \emph{zero-sum}: remembered tokens are tokens no longer
available for reasoning. It is \emph{undifferentiated}: every remembered sentence
is equally ``present'', so there is no distinction between long-term memory and
the current working set. And it is \emph{inescapable}: once something is in the
context it cannot be selectively removed without rewriting the prompt and losing
the KV cache.

\paragraph{Memory in the context may also be behaviourally sticky.} The second
motivation was the hypothesis that content in the context is treated as ``what is
happening now'' rather than as ``history that can be recalled'', and that a memory
living in the weights would therefore be less likely to be carried across a topic
switch. We state this here because Section~\ref{sec:negative} reports that we
\textbf{refuted} it, and because that refutation is more useful to a reader than
another positive result would have been.

\paragraph{What we build.} A frozen backbone plus a side path of $K$ memory slots,
each occupying a contiguous block of a single low-rank adapter's rank dimension
(Section~\ref{sec:method}). A slot is written by gradient steps restricted to its
own block, read implicitly through the forward pass, so the memory content never enters the context, and forgotten by restoring that block to its initial state, which makes forgetting \emph{verifiable}.

\paragraph{Contributions.}
\begin{enumerate}\itemsep2pt
\item A rank-blocked side path whose three operations are independently checkable:
  writes are asserted to touch one slot only, erasure is asserted to return the slot
  to its virgin state bit-for-bit, and reads can be ablated slot-by-slot on a
  byte-identical model. The same property holds at capacity and across a process
  restart (Tables~\ref{tab:t1}, \ref{tab:t4}, \ref{tab:t6}).
\item The finding that read-time \emph{composition} is the binding constraint: the sum
  of the slots is worthless, one selected slot is as good as an oracle, and two
  selected slots already lose most of the benefit. The residual gap equals the routing
  error exactly, on both backbones we test, and the key that drives the selection turns
  out to matter more than the slot mechanism: a model-free lexical overlap lifts 1.7B
  routing from $\tKeyLastTopLarge{}$ to $\tKeyBestTopLarge{}$ (Table~\ref{tab:t2}).
\item A label-free write criterion that reaches full-write recall with
  $\tSelfWrites$ writes instead of $\tAlwaysWrites$ and $\tSelfKL$ nats of backbone
  drift instead of $\tAlwaysKL$, together with a refutation of the likelihood-based
  criterion it replaces (Tables~\ref{tab:t3}, \ref{tab:t3loss}).
\item A scoped diagnostic on real benchmark content in which the \emph{same} router
  delivers the \emph{same} information through the context, graded by a judge
  calibrated against human marks: the parametric arm answers every question correctly
  at zero extra prompt tokens, while context delivery pays
  $\tTokContextTarget{}$--$\tTokContextAll{}$. At equal training budget the bank also
  beats writing those memories into the weights, and it is the only arm that can erase
  one of them afterwards (Table~\ref{tab:t7}).
\item Negative results we report instead of burying: parametric memory does not reduce
  context inertia; a likelihood threshold measures wording; and a retention term that
  looked like a repair of two-slot recall at $n = \tRetentionSmallProbes{}$ does not
  replicate at $n = \tRetentionProbes{}$, although it does cut backbone damage from
  \tRetentionKlBase{} to \tRetentionKlBest{} nats
  (Tables~\ref{tab:t8}, Section~\ref{sec:negative}).
\end{enumerate}

\section{Related work}
\label{sec:related}

\paragraph{Runtime parametric memory.} Titans~\cite{titans} is the clearest
precedent: a neural long-term memory updated at test time by surprise, with a
decay that acts as forgetting. It changes the backbone, whereas we freeze it.
GradMem~\cite{gradmem} writes contextual information into memory tokens by
test-time gradient descent and is the closest work to our setting, but it has no forgetting mechanism, which is the axis we measure most carefully.
TTT-style methods~\cite{ttt,ttte2e} compress context into weights and report
latency wins at long context; \cite{ttte2e} also documents that the naive
per-token cost of full attention at 128K is $2.7\times$ what its method needs.
None of these provide per-item erasure.

\paragraph{Slots and adapters.} ProCL~\cite{procl} organizes LoRA adapters into
structured \emph{program memory} slots retrieved by input-conditioned attention,
motivated by complementary learning systems. Architecturally it is the nearest neighbour to our slot idea. Its setting is task-incremental fine-tuning with task
datasets, it does not forget, and its slots are separate adapters rather than rank
blocks inside one adapter. WISE~\cite{wise} observes an impossibility triangle for
lifelong editing (reliability, generalization, locality) and uses a main/side
memory split, which is the structure we instantiate. AlphaEdit~\cite{alphaedit}
projects edits into the null space of preserved knowledge, a technique we apply
per slot in the anchor term. What none of them offer is a slot that can be written,
read and erased independently at runtime on a frozen backbone; that combination is
what we test, and the erasure half of it is what makes the bit-identical claims
possible.

\paragraph{Fact editing.} ROME/MEMIT-style editors~\cite{rome,memit} write facts
into weights offline and are evaluated mostly on ripple consistency;
\cite{rippleedits} reports that a simple context-editing baseline can top that
metric, a warning that we take seriously (Section~\ref{sec:limits}). Our protocol
therefore carries that baseline explicitly: the context arms of
Table~\ref{tab:t7} receive the same facts by editing the prompt, and they are the
comparison we report against.

\paragraph{Agent memory.} MemGPT/Letta, Mem0, A-MEM, MemoryBank, MemoRAG and
Zep~\cite{memgpt,mem0,amem,memorybank,memoRAG,zep} are the systems practitioners
actually use. Every one of them stores memory in a context layer or an external
store and retrieves it back into the prompt; none performs runtime parameter
updates. That leaves our setting genuinely unoccupied, and it also means there is no established baseline protocol to inherit. Section~\ref{sec:protocol} defines one.

\section{Method}
\label{sec:method}

\subsection{Rank-blocked slots and exact erasure}

The backbone $M_\theta$ is frozen. Each targeted linear projection $W_0$ gets one
low-rank adapter $(A, B)$ with total rank $R = K \cdot r$, partitioned into $K$
contiguous blocks of $r$ rows (for $A$) and columns (for $B$). The effective
weight is
\begin{equation}
W = W_0 + \frac{\alpha}{r} \, B A , \qquad
\Delta = \frac{\alpha}{r}\sum_{k \in \mathrm{active}} B_k A_k .
\end{equation}
Three properties follow, and each is enforced by an assertion.

\emph{Writing isolates.} A write activates a gradient mask that zeroes every rank
block but the target, and builds a \textbf{fresh optimiser} for each write.
The second detail matters: with a reused optimiser, a frozen slot has gradient
zero but non-zero Adam moments, so it drifts anyway; and AdamW's \emph{decoupled}
weight decay shrinks parameters whose gradient is exactly zero, which silently
erodes memories nobody is writing. We therefore build the optimiser inside the
write call with \texttt{weight\_decay=0.0}, and the writer \emph{asserts}
isolation by comparing every non-target slot against a snapshot, raising if one
moved. Both failure modes are covered by tests.

\emph{Erasing is exact.} Each slot keeps a copy $A^{\mathrm{init}}_k$ of its
initial block; \texttt{erase($k$)} restores it and zeroes $B_k$. The result is
bit-identical to a slot that was never written, so ``forgotten'' is an assertable
state, not a small number. Table~\ref{tab:t1} shows the behavioural consequence:
with every slot erased the model's output is \emph{string-identical} to the
prompt-only condition.

\emph{Reading is ablatable.} Because the read path multiplies $A$ and $B$ by a
per-slot mask, MEM\_ON / MEM\_OFF / MEM\_SHUFFLE are produced on a byte-identical model, so the only difference between arms is which slots contribute.

\subsection{Write objective}

The write loss is
\begin{equation}
\mathcal{L} = \underbrace{\mathrm{CE}(v \mid q)}_{\text{target}}
\; + \; \lambda \underbrace{D_{\mathrm{KL}}\!\left(p_{\text{frozen}}^{(t)} \,\|\, p_{\text{slot}}^{(t)}\right)}_{\text{anchor}},
\label{eq:loss}
\end{equation}
where the target term teaches the pair $(q \to v)$ with the \emph{target slot
alone} active, and the anchor term is evaluated on a small set of generic prompts
whose frozen distributions must not move.

Two details of Equation~\ref{eq:loss} were fixed by measurement, not by design.
First, the anchor term is evaluated \textbf{per slot}, not on the sum. Evaluated on
the sum, the individual slots' anchor displacements \emph{cancel}:
per-slot $D_{\mathrm{KL}}$ was $[0.08, 3.89, 1.80, 1.80]$ while the sum's was
$0.09$, i.e. the guard was worth $\lambda \cdot 0.09$ against a target
cross-entropy of $\approx 15$ and was therefore vacuous. Cancellation is a
coincidence of the current slot set and disappears as soon as the read path
selects a single slot, so the guard must hold for each slot on its own. Second,
the anchor is evaluated at \textbf{every position} of the anchor prompts, not only
the last: a final-position-only guard held the single-step divergence at $0.08$
while generation still collapsed into repetition, because autoregressive decoding
compounds a small per-step shift over the whole trajectory. With both fixes the
degenerate repetition disappears and the sum-based read improves from $1/4$ to
$2/4$; both alternatives remain available as an ablation
(\texttt{--anchor-mode sum}).

The anchor term protects the \emph{backbone}: generic prompts must not move. It does
not protect the \emph{bank}. A write is told nothing about the queries the slots
already answer, which is why reads that activate more than one slot collide
(Section~\ref{sec:results}). We therefore also measure an optional \textbf{retention
term} of the same shape,
\begin{equation}
\mathcal{L} = \mathrm{CE}(v \mid q)
\; + \; \lambda \, D_{\mathrm{KL}}\!\left(p_{\text{frozen}}^{(t)} \,\|\,
p_{\text{slot}}^{(t)}\right)
\; + \; \lambda_{\mathrm{ret}} \, D_{\mathrm{KL}}\!\left(p_{\text{bank}}^{(t)}
\,\|\, p_{\text{bank}+k}^{(t)}\right),
\label{eq:retention}
\end{equation}
whose reference $p_{\text{bank}}$ is captured from the bank \emph{as it stands before the write} (previously written slots active, new slot zero-initialised and therefore contributing nothing) on the queries of earlier memories, and whose
current term is evaluated with the new slot added, i.e. under the read condition
where the interference actually appears. Its weight is swept in
Table~\ref{tab:t8}; $\lambda_{\mathrm{ret}} = 0$ recovers
Equation~\ref{eq:loss}.

\subsection{Reading: select, do not sum}
\label{sec:routing}

With every slot active the answers collapse (Table~\ref{tab:t2}), which is
unsurprising once the slots are seen as additive perturbations of the same
weights: each was trained to produce a low-entropy answer, and their sum is a
dominated mixture. The slots are nevertheless perfectly separable: reading with the \emph{wrong} slot yields nothing. The read path should therefore select.
We use the query key, the final hidden state of the query under the frozen
backbone (captured with the memory switched off, so a key does not drift as
memories accumulate), and activate the single best-matching slot by cosine
similarity. No parameters are added and nothing is trained.

\subsection{What deserves to be written}

We compare three criteria, all label-free with respect to ``which item matters'':
\emph{always} (write everything writeable, the reference point), \emph{surprise}
(write iff the pre-write loss exceeds a threshold), and a \emph{self-check}
(ask the frozen model the question; write only if it cannot already answer it).
An explicit-instruction criterion (``remember this'') is included for contrast.
Information gain, the loss difference before and after a tentative write, is the natural fourth criterion. We deliberately do not implement it: it costs a full extra write per candidate, and after Table~\ref{tab:t3} that cost has to be earned.

\subsection{Forgetting and persistence}

Capacity is the integer $K$. A :class:`MemoryStore` (in the implementation, a
dataclass) tracks per-slot utility, recency, norms and an importance flag;
``important'' slots are excluded from eviction, which is what makes an item a
\emph{thought stamp} in the intended sense. Utility decays with a logical clock,
and five eviction policies are implemented: FIFO, LRU, LFU, utility and utility-with-time. The classic baselines are included deliberately, because the
one study with hard numbers on eviction reports that no policy beats LFU by more
than $0.041$ points \cite{semcache}.

A memory that dies with the process is a session cache, so snapshots store the
slot tensors plus the store's bookkeeping. Loading validates slot count, rank,
$\alpha$, the module set and every tensor shape, and refuses on any mismatch: the
dangerous failure is a \emph{partial} load, which would answer from a misaligned
memory while looking healthy.

\section{Evaluation protocol}
\label{sec:protocol}

\paragraph{A fictional benchmark.} Every entity and value in our episodes is
generated from syllables, so no answer can be produced from pre-training
knowledge. Without this, a correct answer is not attributable to the write. For
the same reason, the ``already known'' class used in Section~\ref{sec:results} is
\textbf{verified by asking the model} rather than assumed: of $19$ real-world
facts, the 0.6B model answered only $13$ correctly, getting $2+2 = 2$, the number
of days in a week $=1$, and the largest planet $=$ Mars. Assuming otherwise would
have invalidated that experiment.

\paragraph{Five isolation checks.} P1 \emph{eviction}: the answer must be absent
from the read-time context by exact, normalised and numeric matching, or the item
is dropped from the metric. P2 \emph{counterfactual}: values are randomised per
seed, so answering the value written \emph{this} run is only possible from the
write. P3 \emph{negative control}: never-written entities must not improve.
P4 \emph{ablation}: with the context byte-identical, only the set of active slots
changes. P5 \emph{prompt-only}: strips the benefit of the prompt format. A claim
is reported only if the paired differences over prompt-only have a bootstrap
confidence interval excluding zero.

\paragraph{Paraphrased probes.} Reading reuses a differently worded question for
the same intent, while the memory is written from the original phrasing. With the
write query reused verbatim, top-1 routing scores $100\%$ and the result is a
string-matching artefact rather than a memory result.

\paragraph{Floor screening.} For the inertia experiment we first measure the
no-memory condition and discard any scenario the model answers from its prior.
Only \tScreenKept{} scenarios survived, and every discarded one was the road-side
case, where ``left'' is the model's default answer.

\paragraph{Real content, and how it is graded.} The public-benchmark diagnostic of
Section~\ref{sec:results} uses LongMemEval's oracle split; each memory item is the
benchmark's own $(q,a)$ pair and the evidence turn is located with the benchmark's
\texttt{has\_answer} flags, i.e. retrieval is oracle-indexed by construction. The
comparison arms receive the \emph{same} information through the context, selected by
the \emph{same} query-key router, so the medium is the only variable that changes.
Grading is semantic. The containment measure used on the synthetic experiments rewards reproducing the reference string, which is exactly what a trained write is for, and it therefore cannot compare arms that answer by paraphrase. So the real-content results are graded by a hosted LLM judge, one
candidate per call, calibrated against \tJudgeMergedMarks{} human marks collected on
an anonymised, shuffled sample with the judge's own verdicts hidden until after
marking. The agreement is \tJudgeMergedRate{} (Wilson 95\% interval
[\tJudgeMergedLow{}, \tJudgeMergedHigh{}]) with \tJudgeMergedFalsePositives{} cases of
the judge accepting an answer the human rejected. We report the size of that
calibration and the fact that the judge prompt was revised once, because both bound
what the agreement is worth: the first prompt, on the same rows, agreed on 8 of the
10 rows whose verdict was usable. Two further
protocol points carry weight: wherever every arm answers the \emph{same} items we compare them paired (exact sign test), not by overlapping intervals, and the generation budget is treated as part of the treatment: every arm is generated with the same limit, set so that the longest reference answer fits, since a budget
below that measures truncation instead of memory.

\section{Results}
\label{sec:results}

\paragraph{Writes land, and only where intended.} Table~\ref{tab:t1} is best read
for its controls rather than its accuracy: the write isolation assertion passes
everywhere, the erased condition is string-identical to prompt-only
($\tEraseIdentical$), and the target cross-entropy after a write is
$\tTargetCE$. The low \texttt{mem\_on} row in that table is the \emph{sum-based}
read and is the failure that Section~\ref{sec:routing} addresses; it is not
evidence about writing. The per-slot evidence is in Table~\ref{tab:t2} (oracle
row) and Table~\ref{tab:t6} (a fresh process recalling from disk).

\paragraph{Composition dominates everything else.} Table~\ref{tab:t2} is the
central result. Writing the same memories and changing only which slots
contribute, the sum scores $\tSum{}$, selecting one slot scores $\tTopOne{}$
against an oracle of $\tOracle{}$, and selecting two already falls to
$\tTopTwo{}$. The router's top-1 accuracy is $\tRouterAcc{}$, and the arithmetic
closes: the oracle-minus-router gap equals the routing error exactly, so the
selection rule introduces no loss of its own. The practical consequence is
stronger than ``sparser is better'': a read path for this kind of memory should
select exactly one slot, and a top-$k$ hedge costs most of the benefit.

The shape survives the scale control. On a 1.7B backbone the sum is still
$\tScaleTwoSum{}$, one selected slot is $\tScaleTwoTopOne{}$ against an oracle of
$\tScaleTwoOracle{}$, two slots are $\tScaleTwoTopTwo{}$, and the same arithmetic
closes exactly (oracle minus top-1 equals the routing error, $1 -
\tScaleTwoRouter{}$). What moves with scale is \emph{where} the loss sits: the router
is markedly weaker at 1.7B ($\tRouterAcc{}$ at 0.6B against $\tScaleTwoRouter{}$), so
the gap to the oracle becomes a retrieval problem rather than an interference one.
That makes the router the component to improve next.

\paragraph{The key, not the slot, is what limits that gap.} Because routing accuracy
converts one-for-one into recall, it is worth asking whether the key is simply the wrong
key, and it is: with the same writes and the same probes, changing only the definition
of the retrieval key moves top-1 routing from $\tKeyLastTop{}$ to $\tKeyBestTop{}$ at
0.6B and from $\tKeyLastTopLarge{}$ to $\tKeyBestTopLarge{}$ at 1.7B, and end-to-end
recall moves with it. Nothing is trained; the winning definitions are the ones that keep
more of the query. The best of them is \tKeyBest{}, a bag-of-words TF-IDF overlap
between the query and the stored query, and a 50/50 blend of it with the hidden-state
cosine is close behind. Two variants lose to the incumbent, and both are worth
reporting: keys read from shallower layers (\tKeyShallowTop{} at 0.6B) discard most of
the signal, and re-ranking the two best-routed candidates by the model's own confidence
is \emph{worse} than trusting the key (\tKeyEntropyTop{}), so a slot being sure of
itself is not evidence that it is the right slot. Paired on the same probes, the gain is
significant at 1.7B (\tKeyPairHelpedLarge{} probes helped, \tKeyPairHurtLarge{} hurt,
$p = \tKeyPairPLarge{}$) and absent at 0.6B (\tKeyPairHelpedSmall{} against
\tKeyPairHurtSmall{}, $p = \tKeyPairPSmall{}$): the key matters more as the backbone
grows. One boundary is essential here. These queries come from our synthetic generator,
where content words are distinctive; real benchmark items about a single entity share
almost all of their surface form, so whether a lexical component helps there is a
separate question, and Section~\ref{sec:limits} leaves it open.

\paragraph{Testing a repair for the composition failure.} The boundary is not
retrieval: the correct slot is among the two best-routed ones on
\tRetentionRouteTwo{} of the probes, yet two-slot recall is only
\tRetentionTopTwoBase{}. The interference is manufactured by the writes: each one is trained to produce its own value and told nothing about the queries the bank already answers. We therefore made it explicit. The write objective gains a
\textbf{retention term} (Equation~\ref{eq:retention}): earlier memories' queries must
keep their next-token distributions under the read condition where interference
actually appears (all written slots active), with the reference captured from the
bank as it stands before the write. It reuses the anchor machinery of
Section~\ref{sec:results}, pointed at real memories instead of generic prompts.

\textbf{It does not repair two-slot recall, and we report it alongside the smaller run that said otherwise.} Over \tRetentionProbes{} probes the paired comparison is
\tRetentionHelped{} probes helped against \tRetentionHurt{} hurt (exact sign test
$p = \tRetentionP{}$), and the two-slot rate moves only from \tRetentionTopTwoBase{} to \tRetentionTopTwoBest{}, which is indistinguishable from noise. An earlier, smaller run
of \tRetentionSmallProbes{} probes had looked like a repair (up to
\tRetentionSmallBest{}, \tRetentionSmallHelped{} helped against
\tRetentionSmallHurt{} hurt, $p = \tRetentionSmallP{}$); the effect did not survive
four times the sample, which is why we report both. Two things \emph{do} hold, and
they hold in every run we did. Backbone damage falls by about four fifths, from
\tRetentionKlBase{} to \tRetentionKlBest{} nats. And the condition the term is
actually evaluated on improves one-directionally: reading \emph{every} slot goes from
\tRetentionAllBase{} to \tRetentionAllBest{}, with \tRetentionAllHelped{} probes helped and \tRetentionAllHurt{} hurt ($p = \tRetentionAllP{}$). The magnitude is small, and the direction holds. The honest summary is that teaching a write to protect its
predecessors buys a large reduction in backbone damage and a small gain in the
many-slot condition, while the two-slot repair we hypothesised is unsupported.

\paragraph{Deciding what to remember.} Table~\ref{tab:t3} compares write
criteria on a stream mixing $5$ fictional facts with $13$ verified-known facts.
Writing everything costs $\tAlwaysWrites$ gradient writes, of which
$\tAlwaysWasted$ are spent on facts the model already answers, and moves the
frozen backbone by $\tAlwaysKL$ nats. The self-check criterion spends
$\tSelfWrites$ writes, wastes $\tSelfWasted$, misses $\tSelfMissed$ targets, and moves the backbone by $\tSelfKL$ nats. That is the same recall for
$\tSelfWriteReduction$ fewer writes and $\tSelfKlReductionMin$--$\tSelfKlReductionMax$
less drift across \tSelfSeeds{} seeds (the counts are deterministic; the drift is
not, which is why it is a range). The likelihood-based criterion is the cautionary
tale:
it separates the classes only barely (Table~\ref{tab:t3loss}), because a fixed
phrasing measures whether the model would use \emph{that wording}, not whether it
knows the fact. Asked ``What is the capital of France?'' the model answers
``France's capital is Paris.''; forcing the continuation \texttt{" Paris"} costs
$\tForcedParisKL$ nats against $\tVerboseParisKL$ for the sentence the model
actually produces, even though the answer is right.

\paragraph{Forgetting is exact, and only the dumb policy is bad.}
Table~\ref{tab:t4} streams $12$ memories into $6$ slots, of which the first $3$
are accessed repeatedly. Every eviction restores its slot to the virgin state:
\tErasedTotal{} slots verified bit-identical to never-written, \tErased{} of them
in each policy. The slot flagged important survives every overflow. FIFO keeps only $\tFifoHot$ of the memories that were in use: it evicts what you use and keeps what you never touched. The other four policies keep $\tSmartHot$: the policy matters, but the ``smart'' policies tie, which
independently reproduces the literature's finding that eviction heuristics are
easy to over-credit.

\paragraph{Persistence.} Table~\ref{tab:t6} shows two separate processes. The
read process's side path is virgin before loading ($\tVirgin$), so its
$\tSixOracle$ oracle recall can only come from disk; routed recall is $\tRouted$,
equal to its routing accuracy. Erasing a restored slot still returns it to the
virgin state and the memory does not come back, so forgetting survives a restart. Without that property, a ``forgotten'' memory would be one reboot away from resurrection.

\paragraph{Real content, under oracle indexing.} Table~\ref{tab:t7} reports a
controlled diagnostic on LongMemEval's oracle split~\cite{longmembench}: the 30
single-session questions, history absent throughout, each answer written into its
own slot. The weights arm reproduces \tLongMemOn{} of the answers (verbatim
\tLongMemFull{}), the frozen model answers \tLongMemFrozen{} of them, the erased
condition collapses to \tLongMemOff{} and is bit-identical to the frozen answers,
and the router selects the correct slot every time (\tLongMemRouting). Writing real
answers costs \tLongMemWriteMin{} minutes for \tLongMemSubset{} items and leaves a
mean cross-entropy of \tLongMemCe{}. On a held-out replication over the next 30
questions, containment is \tHoldoutOn{}, \tHoldoutVerdict{} of the pre-registered band (first value $\pm 0.15$). The effect is not an artefact of which questions were drawn.

\paragraph{Why not just keep it in the context?} We measured that too, on the same
questions and with the same router, and graded every arm's answer against the
reference with an LLM judge calibrated against \tJudgeMergedMarks{} human marks: the
parametric arm is correct on \tJudgeWeights{} of the questions and \emph{never loses
one} to any context arm (paired: \tJudgeWins{} won, \tJudgeLosses{} lost).
\texttt{context\_all} reaches \tJudgeContextAll{} while carrying all \tLongMemSubset{} pairs and
paying \tTokContextAll{} extra prompt tokens per question; retrieval into the
context reaches \tJudgeRagOne{} with one item (\tTokRagOne{} tokens) and
\tJudgeRagThree{} with three; a single oracle-selected pair reaches
\tJudgeContextTarget{} for \tTokContextTarget{} tokens; the parametric arm pays
\tTokWeights{}. Two measurement lessons are ours to report, and both make the
apparent gap smaller than a first reading suggests. First, the containment column is
biased: it rewards reproducing the reference string verbatim, which is exactly what
a trained write is for, and it scored correct paraphrases of the same content at
zero. Second, our own first pass at this comparison found the arms
\emph{indistinguishable} (containment \tLongMemOnArchived{} for the parametric arm):
that was an artefact of a 32-token generation budget truncating its answers, which
the judge then reasonably called incomplete. Raising the budget for \emph{every} arm
moved the parametric arm to \tLongMemOn{}. The claim is correspondingly scoped: the
memory items \emph{are} the benchmark's own $(q,a)$ pairs, the evidence turn is
located with the benchmark's own flags (retrieval is oracle-indexed), the write query
is the read query (so \tLongMemRouting{} is an upper bound; paraphrased routing
measured \tRouterAcc{} in Section~\ref{sec:routing}), the calibration rests on
\tJudgeMergedMarks{} marks, and a residual preference for near-verbatim answers cannot
be excluded. A zero observed disagreement is not a zero error rate: with
\tJudgeMergedMarks{} marks the Wilson interval still admits an error rate up to
$1 - \tJudgeMergedLow{}$.

The same comparison on a 1.7B backbone (Table~\ref{tab:t7}) gives \tScaleWeights{} for
the parametric arm against \tScaleContextTarget{}--\tScaleContextAll{} for context
delivery, with the cost column unchanged (\tTokWeights{} against up to
\tScaleTokens{} tokens) and again no question lost by the parametric arm. The margin
does narrow with scale, because a larger backbone uses the context better
(\tScaleFrozen{} of these questions are answered from its prior, against
\tLongMemFrozen{} at 0.6B). So the cost advantage is scale-invariant and the quality advantage is not. That is the honest reading of a two-point curve. This is a
scoped diagnosis, not a LongMemEval result.

\paragraph{Why not simply train the same memories in?} Delivering the facts through the
context is the cheap alternative to a memory bank; training them into the weights is the
expensive one, and a reviewer will ask about it. Three arms, all given the same number of
item-gradient-steps (thirty items, twelve steps each, so the fine-tune arm updates every
parameter on every step and therefore receives \emph{more} compute, not less): full
fine-tuning reaches \tBaseFull{} containment, a single adapter with the same total rank as
the bank reaches \tBaseAdapter{}, and the slot bank reaches \tBaseSlots{}. The bank is
also the gentlest on the frozen backbone, at \tBaseSlotsKl{} nats of all-position KL
against \tBaseFullKl{} for the fine-tune and \tBaseAdapterKl{} for the joint adapter, and
it is the only arm that can forget one memory afterwards: erasing a single slot leaves it
provably virgin (\tBaseErasedVirgin{}), its neighbour still answers
(\tBaseNeighbourKept{}), and the erased question is no longer answered
(\tBaseErasedGone{}). The caveat is that all three arms stop at the bank's budget, so this
is a claim about equal budgets rather than about what fine-tuning could reach given
more.

\section{Negative results}
\label{sec:negative}

\paragraph{Parametric memory does not reduce context inertia.} We pre-registered
the prediction that storing a premise in the weights would drag it into a later,
unrelated topic \emph{less} than leaving it in the context. Table~\ref{tab:t5}
refutes it: inheritance is $\tContext$ for the context arm and $\tParam$ for the
parametric arm, with identical residue. The floor (no premise anywhere) is
$\tFloor$, so the effect is real and not a prior artefact; the control arm, in
which the premise is written and then erased, collapses onto the floor exactly.
In other words the carry-over comes from the model conditioning on memory
\emph{content}, not from that content occupying context tokens. A secondary
observation deserves recording: withholding (the model saying it does not know)
occurs at rate $\tWithholdFloor$ when there is no memory and at
$\tWithholdMemory$ when the memory is present: \emph{having a memory makes the model more confident as well as better informed}.

\paragraph{A likelihood-threshold criterion measures wording.} Described above
and visible in Table~\ref{tab:t3loss}. It is worth stating as a general caution:
``surprise'' is only a proxy for ``unknown'' if the probe is allowed to use the
model's own phrasing. Forcing a canonical answer span makes the proxy measure
verbosity structure instead.

\section{Limitations}
\label{sec:limits}

\begin{itemize}\itemsep2pt
\item \textbf{Scale.} Two of the four headline results are repeated on a 1.7B
  backbone and both survive: the composition shape (Table~\ref{tab:t2}: the sum is
  still worthless and the selection rule still introduces no loss of its own) and the
  real-content diagnostic (Table~\ref{tab:t7}). What does \emph{not} survive is the
  router, whose accuracy falls from $\tRouterAcc{}$ to $\tScaleTwoRouter{}$. That scale effect shifts the bottleneck from interference to retrieval; it does not remove it, although Section~\ref{sec:routing} shows that half of it is a key choice rather than a capacity limit: a model-free lexical key closes the 1.7B gap to $\tKeyBestTopLarge{}$. The inertia refutation and the forgetting experiments remain
  0.6B-only, and 3B and above is untested on this hardware.
\item \textbf{The base model pays.} Writing moves the frozen backbone: after a
  $12$-memory stream the generic prompts are no longer answered verbatim and the
  all-position KL is $\approx 1$ nat. Parametric memory is not free, and the cost
  grows with the number of writes.
\item \textbf{The retention term's headline effect did not replicate.} On
  \tRetentionSmallProbes{} probes it looked like a two-slot repair
  (\tRetentionSmallHelped{} helped to \tRetentionSmallHurt{} hurt,
  $p = \tRetentionSmallP{}$); on \tRetentionProbes{} probes the paired counts are
  \tRetentionHelped{} to \tRetentionHurt{} ($p = \tRetentionP{}$), i.e. nothing. We
  keep both runs in the paper because the shrinkage is the finding. What survives is
  the backbone-damage reduction and a small, one-directional gain in the all-slots
  condition; many-way superposition is still unsolved, and the two-slot limit of
  Section~\ref{sec:routing} therefore still stands as measured.
\item \textbf{Public benchmarks: a scoped diagnosis, not a comparison.} We report
  protocol-controlled synthetic results plus one scoped public-benchmark
  diagnostic (Section~\ref{sec:results}). We did cost the adaptation of
  LongMemEval~\cite{longmembench} on this machine instead of waving at it: its
  oracle split is small enough that a pilot is minutes of compute, so compute is
  \emph{not} the obstacle. Two obstacles are. First, a memory item here is a
  $(q,v)$ pair, while a benchmark turn is a sentence, so a turn$\rightarrow$pair extraction step is required. That is genuine new code, and noisy extraction lands directly in the weights. Second, and more fundamental: the benchmark's hard
  question types (multi-session and temporal reasoning) require combining evidence
  from \emph{several} stored items, whereas our read path deliberately activates
  exactly one slot (Section~\ref{sec:routing} shows two already collapse recall,
  and the retention term of Section~\ref{sec:results} repairs that only partly).
  That is a mechanism boundary, not a compute limit, and it would not move on a
  larger machine. MemoryAgentBench~\cite{mab}, the one public skeleton that
  isolates test-time learning \emph{and} selective forgetting, remains the right
  target once a cross-item read path exists; we do not claim comparability we have
  not measured.
\item \textbf{Two-timescale consolidation is not implemented.} The
  session-fast / cross-session-slow split remains future work, and after the
  inertia refutation its motivation needs restating.
\item \textbf{Metric crudeness.} Base degradation is partly judged by exact
  answer identity, which counts ``$2+2=4$'' as a change from the frozen ``2'';
  the KL numbers are the more informative quantity.
\item \textbf{Single seed for most synthetic stages.} The synthetic experiments
  (T1--T5) were run once each; what we did instead of a full multi-seed matrix was
  to replicate the \emph{public-benchmark diagnostic} on a held-out subset of
  questions (Section~\ref{sec:results}), which is where a lucky draw would have
  mattered most. Seed-level variance for the synthetic results is therefore
  unquantified, and every experiment script takes a seed argument so it can be
  measured.\end{itemize}

\section{Conclusion}

A side path of rank-blocked slots does three things well and one thing only
partly. It writes without disturbing its neighbours, it reads without spending
context, and it forgets in a way that can be verified bit-for-bit, including across a process restart. On real benchmark content, delivering the same information
through the context with the \emph{same} router answers \tJudgeContextTarget{} to
\tJudgeContextAll{} of the questions while paying \tTokContextTarget{} to
\tTokContextAll{} extra prompt tokens per question; the parametric arm answers
\tJudgeWeights{} of them and pays \tTokWeights{}. It does \emph{not} insulate the
model from the previous topic: a memory is behaviourally sticky wherever it is
stored. The practical design rules we would carry forward are that read-time
composition must default to exactly one slot, that the criterion for writing should
ask the model rather than threshold its likelihood, and that capacity is best spent
on durable, exactly erasable slots rather than on trying to keep memories out of the
working set by moving them into weights. The one exception we chased is instructive
in the opposite direction: training each write to protect the memories already in the
bank cuts backbone damage by about four fifths and slightly improves the all-slots
condition, but it does \emph{not} buy back the second slot: an earlier, smaller run suggested it did, and four times the sample refuted it. A measured effect that
shrinks with $n$ is a result worth publishing; we would rather report the refuted
repair than the flattering sample.

\appendix
\section{Reproducibility}

\begin{verbatim}
python -m venv .venv && .venv\Scripts\pip install -e .
python experiments/run_all.py --collect-only      # rebuild tables from reports
python experiments/run_all.py --mode quick       # smoke test, minutes not hours
python experiments/run_all.py --mode full        # re-run every stage from scratch
python paper/build_tables.py                     # regenerate generated_tables.tex
python paper/build_paper_zh.py --compile         # Chinese review copy (a derived aid)
pdflatex main && pdflatex main
\end{verbatim}

Every stage is a separate process writing its own JSON report; a failing stage is
labelled \texttt{FAILED} in the bundle and its numbers render as \texttt{n/a},
never as zero. Table~\ref{tab:status} is that bundle's own summary: which stages are
present in the run that produced this paper. All numbers in this paper come from
\texttt{generated\_tables.tex}, which is generated from \texttt{runs/summary.json}.

Two boundaries are worth stating precisely rather than implying full
reproducibility. First, the T7-f grading phase calls a hosted model: it needs a
chat-completions credential and spends about an hour at the provider's rate limit,
so the reported audit (Table~\ref{tab:t7}) is re-rendered from the stored report
rather than re-derived by \texttt{--collect-only}. The human marks the judge was
calibrated against \emph{are} in the repository
(\texttt{bench/judge\_calibration\_marks.json}), so the calibration can be checked
without an API key. Second, the per-stage reports live in \texttt{runs/}, which is
not pushed (the repository is source-only), so re-deriving the tables means re-running the stages, starting with \texttt{--mode quick}, since a full run is hours.

\begin{thebibliography}{99}\small
\bibitem{titans} A. Behrouz, P. Zhong, V. Mirrokni. Titans: Learning to Memorize
  at Test Time. arXiv:2501.00663.
\bibitem{gradmem} GradMem: Learning to Memorize by Test-Time Gradient Descent on
  Memory Tokens. arXiv:2603.13875 (ICML 2026).
\bibitem{ttt} Y. Sun et al. Learning to (Learn at Test Time): RNNs with
  Expressive Hidden States. arXiv:2407.04620.
\bibitem{ttte2e} Test-Time Training End-to-End. arXiv:2512.23675.
\bibitem{procl} H. Le, S. Venkatesh. Continual Fine-Tuning of Large Language
  Models via Program Memory. arXiv:2605.13162.
\bibitem{wise} WISE: Rethinking the Knowledge Memory for Lifelong Model Editing.
  arXiv:2405.14768 (NeurIPS 2024).
\bibitem{alphaedit} AlphaEdit: Null-Space Constrained Knowledge Editing for
  Language Models. arXiv:2410.02355 (ICLR 2025).
\bibitem{rome} ROME: Locating and Editing Factual Associations in GPT.
  arXiv:2202.05262.
\bibitem{memit} MEMIT: Mass-Editing Memory in a Transformer. arXiv:2210.07229.
\bibitem{rippleedits} RippleEdits: Evaluating the ripple effects of knowledge
  editing. TACL 2024, arXiv:2307.12976.
\bibitem{memgpt} MemGPT: Towards LLMs as Operating Systems. arXiv:2310.08560.
\bibitem{mem0} Mem0: Building Production-Ready AI Agents with Scalable Long-Term
  Memory. arXiv:2504.19413.
\bibitem{amem} A-MEM: Agentic Memory for LLM Agents. arXiv:2502.12110.
\bibitem{memorybank} MemoryBank: Enhancing Large Language Models with Long-Term
  Memory. arXiv:2305.10250.
\bibitem{memoRAG} MemoRAG: Moving towards Next-Gen RAG. arXiv:2409.05591.
\bibitem{zep} Zep: A Temporal Knowledge Graph Architecture for Agent Memory.
  arXiv:2501.13956.
\bibitem{semcache} On eviction policies for semantic caches. arXiv:2608.20280.
\bibitem{mab} MemoryAgentBench: Evaluating Memory in LLM Agents.
  arXiv:2507.05257.
\bibitem{longmembench} LongMemEval: Benchmarking Chat Assistants on Long-Term
  Interactive Memory. arXiv:2410.10813 (ICLR 2025).
\end{thebibliography}

\end{document}
